\documentclass{situationist-preprint}

\title{The Look Without Conflict}
\msid{PHE-09-a}
\msfamily{PHE-09}
\msdiscipline{Phenomenology}
\mstype{Principal manuscript}
\mskeywords{Sartre; the look; being-for-others; objectification; shame; conflict; Beauvoir; Fanon; recognition; Honneth}

\begin{document}
\maketitle

\begin{abstract}
In \emph{Being and Nothingness} Sartre argues that the other is originally given to me in the look, through which I experience myself as an object, and he concludes that conflict is the original meaning of being-for-others. We argue that the analysis of the look establishes objectification but not hostility. We separate three claims: that the look makes me an object for another; that being made an object is an alienation of my possibilities; and that the relation between looker and looked-at is therefore a struggle in which each must objectify the other or be objectified. The first claim is sound. The second holds only in a qualified form, since what the look takes from me is the monopoly of the meaning of my situation, not my possibilities themselves. The third does not follow: it requires the further premise that every objectification is a degradation, and that premise is true only under social conditions in which the look carries a hierarchy. Beauvoir and Fanon describe such conditions; Honneth's account of recognition explains why they are contingent. Sartre's own later work, which ties conflict to scarcity, points in the same direction.
\end{abstract}

\section{The problem}

Sartre's analysis of the look is among the most influential passages in twentieth-century philosophy of other minds \citep[Part~III, ch.~1, \S4]{Sartre1943}. A man looks through a keyhole, absorbed in what he sees; he is nothing but his project of seeing, and is not given to himself as anything. He hears footsteps and realizes that he is being looked at. In that moment he is suddenly something: a voyeur, an object in the world of another, with a nature he did not choose and cannot see from within. Shame is the experience of this: shame is shame of oneself before the other \citep[Part~III, ch.~1]{Sartre1943}. Sartre argues that this experience gives the other's existence with a certainty the argument from analogy cannot provide, and that it reveals a dimension of my being, my being-for-others, that I am but cannot know as I know my own projects.

From this Sartre draws a darker conclusion. In Part~III, chapter~3, he examines concrete relations with others, love, language, masochism, indifference, desire, sadism and hate, and finds that each fails: I try either to assimilate the other's freedom or to objectify it, and both attempts founder. He concludes that conflict is the original meaning of being-for-others, and that the \enquote{we} of shared experience is a psychological rather than an ontological structure \citep[Part~III, ch.~3]{Sartre1943}. \emph{No Exit} puts the conclusion in a slogan: hell is other people \citep{Sartre1944}.

Critics have resisted the conclusion from several directions. Merleau-Ponty argued that the other's look transforms me into an object only if both of us withdraw into our thinking nature and make ourselves an inhuman gaze; communication, not objectification, is the primary relation \citep{MerleauPonty1945}. Beauvoir argued that mutual recognition of freedoms is possible and is the condition of an ethics \citep{Beauvoir1947}. \citet{Theunissen1965} argued that Sartre's analysis inherits the transcendental philosophy's inability to reach a genuinely dialogical relation. \citet{Zahavi2014} has defended the look's role in constituting shame while rejecting Sartre's conflictual conclusion. We argue that the analysis of the look establishes objectification and not hostility, and that conflict is a contingent feature of looks exchanged under particular social conditions.

\section{Framework}

\begin{definition}[Objectification]\label{d:obj}
I am \emph{objectified} by another if, in being looked at, I experience myself as having a determinate character in the other's world, a character I do not confer on myself and cannot apprehend from within my own projects.
\end{definition}

\begin{definition}[Hostility]\label{d:host}
A relation between two subjects is \emph{hostile} if each, in that relation, has its possibilities diminished by the other and can recover them only by diminishing the other's.
\end{definition}

\begin{definition}[Degradation]\label{d:degr}
An objectification is a \emph{degradation} if the character conferred is one that ranks the looked-at below the looker in a hierarchy of worth, so that being so seen reduces the standing of the looked-at in the shared world.
\end{definition}

\begin{principle}[Possibilities are not conferred meanings]\label{p:poss}
What I can do in a situation is not the same as the meaning my doing it has for others. A look can change the second without changing the first.
\end{principle}

Principle~\ref{p:poss} is a distinction Sartre draws himself, in other terms: the other alienates my possibilities by seeing them as possibilities of an object, \enquote{dead-possibilities}, but does not remove the transcendence by which I project them \citep[Part~III, ch.~1]{Sartre1943}. The principle states the distinction as one between the content of my possibilities and their significance in another's world.

\section{The argument}

\begin{lemma}\label{l:obj}
The look objectifies.
\end{lemma}

\begin{proof}
This is Sartre's phenomenological result, and we accept it. In the keyhole case, the man experiences himself as a voyeur: a determinate character, which he did not give himself (he was simply looking) and which he cannot see from within his project of looking. By Definition~\ref{d:obj}, he is objectified.
\end{proof}

\begin{lemma}\label{l:notalien}
Objectification does not by itself diminish the looked-at's possibilities. It removes only his monopoly over the meaning of his situation.
\end{lemma}

\begin{proof}
After the footsteps, the man can still do what he could do before: go on looking, stand up, walk away, explain. What has changed is that each of these acts now has a meaning in another's world, a meaning he does not control. By Principle~\ref{p:poss}, a change in the meaning one's acts have for others is not a change in what one can do. So objectification takes from the looked-at the position of sole author of the meaning of his situation, but it leaves his possibilities intact.
\end{proof}

\begin{theorem}\label{t:noconf}
The look is not essentially hostile. It is hostile only if the objectification it effects is a degradation.
\end{theorem}

\begin{proof}
By Definition~\ref{d:host}, a hostile relation diminishes each party's possibilities. By Lemma~\ref{l:notalien}, objectification as such does not. So the hostility of a look must come from something besides objectification. Consider what would turn the loss of monopoly into a diminution of possibilities. If the character conferred ranks me below the looker (Definition~\ref{d:degr}), then being so seen does diminish what I can do in the shared world: my word counts for less, my movements are watched for threat, my acts are interpreted in advance. In that case the only way to recover my possibilities is to contest the ranking, which means imposing a counter-look; the relation becomes a struggle in the sense of Definition~\ref{d:host}. If the character conferred is not a degradation, for example if I am seen as attentive, or tired, or simply as someone there, nothing in the look diminishes my possibilities, and nothing calls for a counter-look. Hence the look is hostile only if its objectification degrades, and since not every objectification degrades, the look is not essentially hostile.
\end{proof}

Sartre's chapter on concrete relations assumes, at each turn, that to be an object for the other is to be at the other's disposal, so that my freedom can be recovered only by reducing the other to an object. The keyhole case is chosen to support this, since the character conferred, voyeur, is shameful. But shame is one response among several to being seen. Pride is another, which Sartre also discusses; so is the ordinary comfort of being seen by a friend, which he does not. The look of the friend objectifies me in the sense of Lemma~\ref{l:obj}, since I am seen as having a character I cannot see from within, but it does not degrade, and it is not answered by a counter-look.

\section{The social conditions of conflict}

If conflict depends on degradation, the question becomes when looks degrade. Theorem~\ref{t:noconf} implies that the answer lies in the hierarchies that the look carries, and these are social.

\citet{Beauvoir1949} describes the look under which woman is constituted as the Other: she is seen as the inessential correlate of the essential subject, and the meanings conferred on her body and conduct are fixed in advance by that relation. The objectification Sartre describes as symmetrical becomes, under this condition, systematically one-sided. Beauvoir's diagnosis is that this asymmetry is a social and historical formation, not an ontological necessity, and her ethics of ambiguity holds that the freedom of each requires the freedom of others, so that reciprocal recognition is the goal of a liberating practice \citep{Beauvoir1947}.

\citet{Fanon1952} describes the experience of being seen as Black in a white world: the look of a child on a train who points him out to his mother by his color, and the sudden fixing of his body under a historico-racial schema that precedes and overrides his own bodily schema. Fanon's description is a case of Sartrean objectification, and he engages Sartre directly. But it shows that what makes this look crushing is not that it objectifies, which every look does, but that it imposes a degrading character fixed by colonial history. \citet{Gordon1995} has developed Fanon's analysis into an account of antiblack racism as institutionalized bad faith. In both Beauvoir's and Fanon's cases, the conflict is real and the look is its medium, but its source is a social hierarchy that the look transmits.

\citet{Honneth1992} provides the theory that explains why these conditions are contingent. On his account, recognition takes three forms, love, respect and esteem, and social conflict arises from experiences of disrespect, in which one is denied a form of recognition one is due. Struggles for recognition are not the original meaning of intersubjectivity but responses to its distortion. Honneth's account and Theorem~\ref{t:noconf} agree: conflict arises where looks degrade, and looks degrade where social relations deny recognition. Where they do not, being seen is a condition of self-relation, as PID-12 argues in holding that the self is partly constituted by relations of recognition.

Sartre's later work moves toward this view. In the \emph{Critique of Dialectical Reason} he grounds the antagonism of human relations in scarcity: each person is, in a world of scarcity, a possible threat to the other's survival, and it is this material condition that turns reciprocity into conflict \citep{Sartre1960}. Scarcity is a contingent condition, and conflict that depends on it is contingent too.

\section{Objections and replies}

\begin{objection}[Every look degrades]\label{o:every}
To be seen as anything determinate is to be reduced from a free transcendence to a thing. That is already a degradation, whatever the character conferred.
\end{objection}

\begin{reply}
This identifies objectification with degradation by definition, and thereby makes Sartre's conclusion analytic. But it is not true to the phenomenology. Being seen as one's friend's friend is being seen as something determinate, and it does not reduce one's standing in the shared world. The objection must hold that every determination is a reduction because freedom is the absence of determination; but on Sartre's own view freedom is always exercised in a situation that has a determinate character, and being seen is one more feature of the situation in which I act.
\end{reply}

\begin{objection}[The failure of love]\label{o:love}
Sartre's analysis of love shows that even the most benevolent relation fails: the lover wants to be the whole world for the beloved's freedom, which is impossible.
\end{objection}

\begin{reply}
The analysis shows that love conceived as possession of another's freedom fails. It does not show that love as reciprocal recognition fails, and Sartre's own posthumous \emph{Notebooks for an Ethics} explore an authentic love in which each wills the other's freedom \citep{Sartre1983}. The failure of the possessive project is evidence for Theorem~\ref{t:noconf}, not against it: the project fails because it treats the other's look as something to be controlled, and that treatment presupposes that being seen is a threat.
\end{reply}

\begin{objection}[Contingency is no consolation]\label{o:cont}
If the social conditions of degradation are universal in human history, the claim that conflict is contingent is empty.
\end{objection}

\begin{reply}
The claim is about the structure of the look, not about the frequency of conflict. Even if every society had imposed degrading hierarchies, that would show a pervasive social fact, not an ontological one, and the difference matters: an ontological conflict could not be remedied, whereas a social one can be addressed by changing the conditions that produce it. Beauvoir, Fanon and Honneth each draw the practical conclusion that the remedy is political. Moreover, the conditions are not universal across relations even within a society; friendship and solidarity exist under the same social order as degradation.
\end{reply}

\section{Conclusion}

Sartre's analysis of the look is a correct account of objectification: being seen, I acquire a character in another's world that I did not confer and cannot see from within. It does not establish hostility. Objectification removes my monopoly over the meaning of my situation but not my possibilities, and it becomes a struggle only when the character conferred is a degradation. Whether looks degrade depends on social hierarchies, which Beauvoir and Fanon describe and Honneth's theory of recognition explains. Conflict is not the original meaning of being-for-others, but a contingent form it takes under particular conditions.

\begin{thebibliography}{99}\small
\bibitem[Beauvoir(1947)]{Beauvoir1947} Beauvoir, S. de (1947). \emph{Pour une morale de l'ambigu\"it\'e}. Gallimard. Trans. B. Frechtman, \emph{The Ethics of Ambiguity}, Philosophical Library, 1948.
\bibitem[Beauvoir(1949)]{Beauvoir1949} Beauvoir, S. de (1949). \emph{Le deuxi\`eme sexe}. Gallimard. Trans. C. Borde and S. Malovany-Chevallier, \emph{The Second Sex}, Knopf, 2010.
\bibitem[Fanon(1952)]{Fanon1952} Fanon, F. (1952). \emph{Peau noire, masques blancs}. \'Editions du Seuil. Trans. R. Philcox, \emph{Black Skin, White Masks}, Grove Press, 2008.
\bibitem[Gordon(1995)]{Gordon1995} Gordon, L.~R. (1995). \emph{Bad Faith and Antiblack Racism}. Humanities Press.
\bibitem[Honneth(1992)]{Honneth1992} Honneth, A. (1992). \emph{Kampf um Anerkennung}. Suhrkamp. Trans. J. Anderson, \emph{The Struggle for Recognition}, Polity, 1995.
\bibitem[Merleau-Ponty(1945)]{MerleauPonty1945} Merleau-Ponty, M. (1945). \emph{Ph\'enom\'enologie de la perception}. Gallimard. Trans. D.~A. Landes, \emph{Phenomenology of Perception}, Routledge, 2012.
\bibitem[Sartre(1943)]{Sartre1943} Sartre, J.-P. (1943). \emph{L'\^etre et le n\'eant}. Gallimard. Trans. H.~E. Barnes, \emph{Being and Nothingness}, Philosophical Library, 1956.
\bibitem[Sartre(1944)]{Sartre1944} Sartre, J.-P. (1944). \emph{Huis clos}. Gallimard.
\bibitem[Sartre(1960)]{Sartre1960} Sartre, J.-P. (1960). \emph{Critique de la raison dialectique}. Gallimard.
\bibitem[Sartre(1983)]{Sartre1983} Sartre, J.-P. (1983). \emph{Cahiers pour une morale}. Gallimard. Trans. D. Pellauer, \emph{Notebooks for an Ethics}, University of Chicago Press, 1992.
\bibitem[Theunissen(1965)]{Theunissen1965} Theunissen, M. (1965). \emph{Der Andere: Studien zur Sozialontologie der Gegenwart}. de Gruyter.
\bibitem[Zahavi(2014)]{Zahavi2014} Zahavi, D. (2014). \emph{Self and Other: Exploring Subjectivity, Empathy, and Shame}. Oxford University Press.
\end{thebibliography}

\end{document}
